\subsection{Analisis Validasi Aktor dan Operasi Tulis}
\label{subsec:analisis-validasi-aktor-operasi-tulis}

\textit{Caveats} pada token dapat menyatakan bahwa suatu aktor memperoleh hak baca atau tulis terhadap kategori data tertentu. Namun, dalam konteks RME, hak tulis tidak hanya bergantung pada token yang dibawa, tetapi juga pada kewenangan profesional aktor. Operasi tulis memiliki dampak langsung terhadap integritas data klinis karena data yang ditulis dapat menjadi dasar pengambilan keputusan medis berikutnya.

Pada DecMed awal, \textit{role} \texttt{MedicalPersonnel} digunakan untuk merepresentasikan tenaga medis secara umum. \textit{Role} tersebut belum cukup untuk membedakan dokter, perawat, petugas laboratorium, dan apoteker. Padahal, keempat aktor tersebut memiliki kewenangan tulis yang berbeda. Perawat dapat menulis data asesmen awal, dokter dapat menulis diagnosis dan terapi, petugas laboratorium dapat menulis hasil pemeriksaan laboratorium, sedangkan apoteker dapat menulis data terkait peresepan dan dispensing.

Oleh karena itu, diperlukan validasi tambahan terhadap \textit{subrole} tenaga medis, terutama untuk operasi tulis. Token dapat digunakan untuk menyatakan cakupan akses yang diberikan, tetapi \textit{Server} PRE tetap perlu memastikan bahwa aktor aktif memiliki \textit{subrole} yang sesuai dengan \textit{dataset} atau fungsi klinis yang akan ditulis. Validasi ini mencegah kondisi ketika aktor membawa token dengan izin tulis, tetapi secara kewenangan profesional tidak berhak memperbarui segmen data tertentu.

Berdasarkan analisis tersebut, validasi subrole tenaga medis terhadap dataset menjadi dasar bagi \texttt{S-TA-05}. Solusi ini melengkapi penegakan otoritas pada \textit{Server} PRE dengan memastikan bahwa operasi tulis tidak hanya sah berdasarkan token, tetapi juga sesuai dengan kewenangan klinis aktor aktif.
